\section{Introduction}
\label{sec:intro}

The remarkable success of deep neural networks has led to an abundance of specialized models, each fine-tuned from a shared pre-trained base to excel at specific downstream tasks \cite{langley00, DudaHart2nd}. Consolidating these diverse, task-specific capabilities into a single, unified multi-task model without retraining the entire network remains a central ambition of modern machine learning. To this end, weight-space model merging has emerged as an incredibly cost-effective paradigm \cite{mitchell80}. Rather than retraining or maintaining expensive multi-branch architectures, weight-space merging directly interpolates the parameters of fine-tuned expert models.

The simplest approach is \emph{static model merging} (e.g., task arithmetic), which linearly combines task vectors (the differences between expert parameters and base parameters) using fixed, uniform coefficients. However, static merging is inherently limited; a single set of static weights must represent all possible inputs, inevitably leading to representation interference and the loss of fine-grained task-specific capabilities. 

To overcome this fundamental limitation, recent works have proposed \emph{dynamic weight-space merging} (or dynamic routing), where input-dependent routing coefficients are predicted sample-by-sample on-the-fly. For an input sample $x$, a parametric routing module (e.g., a linear-Softmax layer) processes intermediate features to predict a probability distribution over the available experts. The final model parameters are then dynamically assembled as a convex combination of the experts' weights.

\begin{figure}[t]
  \centering
  \begin{picture}(220, 110)
    \put(10, 90){\framebox(40,15){Input $x$}}
    \put(30, 90){\vector(0,-1){20}}
    \put(10, 55){\framebox(45,15){Backbone}}
    \put(32.5, 55){\vector(0,-1){20}}
    \put(10, 20){\framebox(45,15){Features $z(x)$}}
    \put(55, 27.5){\vector(1,0){30}}
    
    \put(85, 20){\framebox(50,15){Proj. $\psi(x)$}}
    \put(110, 35){\vector(0,1){20}}
    \put(85, 55){\framebox(50,15){Router $W, B$}}
    \put(110, 70){\vector(0,1){20}}
    \put(80, 90){\framebox(60,15){Coefficients $\alpha(x)$}}
    
    \put(140, 97.5){\vector(1,0){25}}
    \put(165, 90){\framebox(50,15){Merged $W(x)$}}
    \put(190, 90){\vector(0,-1){40}}
    \put(165, 35){\framebox(50,15){Inference}}
  \end{picture}
  \caption{Dynamic weight-space routing pipeline. Peninsula features $z(x)$ are projected to a unit-state $\psi(x)$, passed to a linear-Softmax router to output coefficients $\alpha(x)$, which dynamically merge expert parameters on-the-fly.}
  \label{fig:pipeline}
\end{figure}

While dynamic routing offers substantial performance gains, it introduces a severe generalization challenge under extreme data scarcity. In many practical scenarios, calibration datasets are exceptionally small (e.g., a batch size of $B_{\text{cal}} \le 64$ samples across all tasks). In this sparse-data regime, parametric routers suffer from severe, low-data overfitting. During calibration, the router's weights grow unchecked to minimize the cross-entropy loss, creating a massive generalization gap. At test time, especially under representation noise or domain shifts, unregularized routing collapses catastrophicly, failing to outperform simple static ensembling and degrading expert capabilities on complex or out-of-distribution (OOD) tasks.

To prevent this collapse, existing methodologies employ ad-hoc, heuristic regularizations. For instance, Task-Space Anchor Regularization (TSAR) anchors the routing weights to predefined task-space centroids, while VR-Router enforces a variance-suppressing penalty on predicted task coefficients. Although these methods provide some empirical relief, they are fundamentally heuristic and lack theoretical justifications explaining \emph{why} they work or \emph{how} they relate to the underlying parameter spaces of the merged models. Specifically, they treat all experts identically, applying isotropic penalties that are blind to the geometry of the expert parameter spaces.

In this work, we reject heuristic ensembling and approach dynamic routing regularization through the lens of first-principles statistical learning theory. We ask a fundamental question: \emph{How does the parameter-space geometry of the expert models dictate the generalization complexity of the dynamically merged model?}

To answer this, we formalize the dynamic weight-space routing problem and derive the first formal upper bound on the Rademacher complexity of a dynamically merged model class. Our derivation reveals that the generalization gap of a dynamically merged model is bounded by a weighted sum of the routing weight norms, where the weights are precisely the parameter-space distances (Frobenius or Spectral norms) of the task-specific expert models from the base model. This formal result demonstrates a critical flaw in uniform $L_2$ decay: since different experts deviate from the base model by vastly different magnitudes, uniform regularization is structurally suboptimal. Experts with large task-vector norms are highly sensitive; a small change in their routing coefficients causes a massive shift in the parameter space of the merged model, inflating the Rademacher complexity.

Guided by this insight, we propose \textbf{Spectral and Rademacher-guided Routing Regularization (SR3)}. Under SR3, we scale the regularization penalty of each task's routing weights proportionally to the squared task-vector norm (either Frobenius or Spectral). This forces the router to be highly conservative when activating distant, high-complexity experts under representation uncertainty, while permitting flexible routing for nearby, low-complexity experts.

Our key contributions are:
\begin{enumerate}
    \item \textbf{Generalization Theory for Merged Models:} We derive the first Rademacher complexity generalization bound for dynamically merged models, mathematically formalizing the relationship between parameter-space geometry and representation-space complexity.
    \item \textbf{Provably Optimal Regularizer (SR3):} We prove that scaling the routing parameter weight decay proportionally to the squared task-vector norm is the theoretically optimal strategy to minimize the generalization bound, and we propose two variants: Frobenius (SR3-F) and Spectral/Operator (SR3-S).
    \item \textbf{Empirical Validation under Fair Conditions:} We evaluate our methods on an unperturbed continuous weight-merging simulator, removing manual baseline scaling multipliers. We show that our mathematically derived SR3 variants achieve highly competitive performance compared to state-of-the-art heuristics, but with the unique advantage of being theoretically justified.
    \item \textbf{Concentration of Measure Analysis:} We provide a deep theoretical analysis of the relationships between Spectral and Frobenius norms in high dimensions. We demonstrate that for random Gaussian task vectors, the concentration of measure forces the spectral norm to be a constant fraction of the Frobenius norm, and we explain how structured real-world task vectors deviate from this behavior to differentiate the two variants.
\end{enumerate}